\documentclass[a4paper,12pt,oneside,fontset=none]{ctexbook}

% ========================================
%  ↓↓↓ 讲义信息设置（修改讲义时只需改这里） ↓↓↓
% ========================================

\newcommand{\LectureNotesTitle}{考研数学笔记}
\newcommand{\AuthorInfo}{@sunlight}

\newcommand{\MakeTitlePage}{%
    \begin{titlepage}
        \centering
        \vspace*{5cm}
        
        {\fontsize{48pt}{60pt}\selectfont\heavykaishu \color{LogicColor} 考研数学笔记\par}
        
        \vfill
        
        {\LARGE \kaishu 
        考研数学笔记
        \\[0.8em]
        \Large
        sunlight
        \\
        }
    
        \vspace{2.3cm}
        
        {\Large \kaishu
        学\ 海\ 悠\ 悠\quad 研\ 途\ 漫\ 漫
        \\[0.3em]
        {\Large 以此薄卷赠诸生共勉}
        \\[2.5em]
        
        %二\ 〇\ 二\ 六\ 年\ 三\ 月 
        }
        
        \vspace{2cm}
    \end{titlepage}%
}

% ========================================
%  ↑↑↑ 讲义信息设置（修改讲义时只需改这里） ↑↑↑
% ========================================

% ======================================
%  ↓↓↓ 模板核心设置（可根据实际需求修改） ↓↓↓
% ======================================

\usepackage[a4paper, left=2.5cm, right=2.5cm, top=2.5cm, bottom=2.5cm]{geometry}
\usepackage{fancyhdr}
\usepackage{amsmath, amssymb, amsthm, mathtools}
\usepackage{mathpazo}
\usepackage[most]{tcolorbox}
\usepackage{xcolor}
\usepackage{newunicodechar}
\usepackage{etoolbox}

\usepackage{esint}
\usepackage{multirow}
\usepackage{diagbox}
\usepackage{needspace}
\usepackage{tikz}
\usepackage{bm}
\usetikzlibrary{arrows.meta}
\usepackage{makecell}
\usepackage{pgfplots}
\pgfplotsset{compat=1.18}
\usepackage{capt-of}
\usepackage{graphicx,caption,subcaption}
\renewcommand{\thesubfigure}{\Alph{subfigure}}
\usepackage{tikz-3dplot}
\tdplotsetmaincoords{70}{110}
\usepackage{booktabs}
\usetikzlibrary{patterns}
\usepgfplotslibrary{fillbetween}


\pgfplotsset{
    % 让x轴标签出现在箭头最右侧
    every axis x label/.style={
        at={(ticklabel* cs:1)},
        anchor=west,
    },
    % 让y轴标签出现在箭头最上方
    every axis y label/.style={
        at={(ticklabel* cs:1)},
        anchor=south,
    }
}


\setlength{\headheight}{26pt}

\newunicodechar{。}{.\ \ignorespaces}
\newunicodechar{，}{,\ \ignorespaces}
\newunicodechar{：}{:\ \ignorespaces}
\newunicodechar{；}{;\ \ignorespaces}
\newunicodechar{？}{?\ \ignorespaces}
\newunicodechar{！}{!\ \ignorespaces}
\newunicodechar{（}{\textup{(}\ignorespaces}
\newunicodechar{）}{\textup{)}\ignorespaces}

\newcommand{\R}{\mathbb{R}}
\newcommand{\sep}{\text{ · \ }}
\renewcommand{\ge}{\geqslant}
\renewcommand{\geq}{\geqslant}
\renewcommand{\le}{\leqslant}
\renewcommand{\leq}{\leqslant}


\AddToHook{env/pmatrix/before}{\linespread{1.2}\selectfont}
\AddToHook{env/bmatrix/before}{\linespread{1.2}\selectfont}
\AddToHook{env/vmatrix/before}{\linespread{1.2}\selectfont}
\AddToHook{env/dcases/before}{\linespread{1.25}\selectfont}
\AddToHook{env/enumerate/before}{\linespread{1.4}\selectfont}
\AddToHook{env/enumerate/before}
{\vrule width 0em height 0cm depth 0.5cm}

\AddToHook{env/pmatrix/after}
{\vrule width 0em height 0cm depth 1.5cm}
\AddToHook{env/bmatrix/after}
{\vrule width 0em height 0em depth 1.5em}
\AddToHook{env/vmatrix/after}
{\vrule width 0em height 0em depth 1.5em}
\AddToHook{env/dcases/after}
{\vrule width 0em height 2em depth 2em}


\let\olddfrac\dfrac
\renewcommand{\dfrac}[2]{
    \olddfrac{#1\vrule width 0pt height 1.15em depth 0em}{#2\vrule width 0pt height 0em depth 0.6em}
}

\let\oldlim\lim
\renewcommand{\lim}{\displaystyle\oldlim}
\let\oldint\int
\renewcommand{\int}{\displaystyle{\vrule width 0pt height 0em depth 0em}\oldint}
\let\oldiint\iint
\renewcommand{\iint}{\displaystyle{\vrule width 0pt height 0em depth 0em}\oldiint}
\let\oldiiint\iiint
\renewcommand{\iiint}{\displaystyle{\vrule width 0pt height 0em depth 0em}\oldiiint}
\let\oldoint\oint
\renewcommand{\oint}{\displaystyle{\vrule width 0pt height 1.5em depth 1.5em}\oldoint}
\let\oldoiint\oiint
\renewcommand{\oiint}{\displaystyle{\vrule width 0pt height 1.5em depth 1.5em}\oldoiint}
\let\oldsum\sum
\renewcommand{\sum}{\displaystyle\oldsum}


\setmainfont{TeX Gyre Pagella}

% 显式绑定系统字体
\setCJKmainfont[
    Path = C:/Users/sunlight/AppData/Local/Microsoft/Windows/Fonts/,
    Extension = .otf,
    UprightFont = *-Regular,
    BoldFont = *-Bold,
    ItalicFont = *-Regular,
    BoldItalicFont = *-Bold,
]{NotoSerifCJKsc}
\setCJKsansfont[
    Path = C:/WINDOWS/Fonts/,
    Extension = .ttf,
]{simhei}
\setCJKmonofont[
    Path = C:/WINDOWS/Fonts/,
    Extension = .ttf,
]{STKAITI}

\newCJKfontfamily\kaishu{KaiTi}
\newCJKfontfamily\heavykaishu{KaiTi}[AutoFakeBold = 2]

\definecolor{LogicColor}{RGB}{200, 80, 110}
\definecolor{LogicBg}{RGB}{255, 240, 245}
\definecolor{DefColor}{RGB}{210, 105, 30}
\definecolor{DefBg}{RGB}{255, 250, 240}
\definecolor{ExColor}{RGB}{70, 130, 180}
\definecolor{ExBg}{RGB}{240, 248, 255}
\definecolor{NoteColor}{RGB}{34, 139, 34}
\definecolor{MainText}{RGB}{20, 20, 20}
\color{MainText}

% --- 新增题型辅助命令 ---
\newcommand{\fillin}[1][3cm]{\,\underline{\hspace{#1}}\,}
\newcommand{\selectbracket}{\unskip\nolinebreak\quad \raisebox{0.7pt}[0pt][0pt]{\mbox{$(\qquad)$}}\par}
% -----------------------

\ctexset{
    chapter = {
        format = \centering,
        nameformat = \fontsize{40pt}{50pt}\selectfont\heavykaishu\color{LogicColor},
        titleformat = \fontsize{40pt}{50pt}\selectfont\heavykaishu\color{LogicColor},
        aftername = \hspace{1.8em},
        beforeskip = 10pt,
        afterskip = 55pt,
    },
    section = {
        format = \centering\LARGE\bfseries\color{MainText},
        beforeskip = 24pt,
        afterskip = 27pt,
    }
}

\tcbset{
    commonbox/.style={
        enhanced,
        fonttitle=\bfseries\large,
        attach boxed title to top left={xshift=5mm, yshift*=-\tcboxedtitleheight/2},
        boxed title style={
            frame hidden,
            opacityback=1,
            rounded corners,
            arc=4pt,
            top=1mm, bottom=1mm, left=2mm, right=2mm
        },
        boxrule=0pt, frame hidden,
        sharp corners=downhill, rounded corners, arc=10pt,
        drop fuzzy shadow=black!15!white,
        left=8mm, right=5mm, top=3mm, bottom=5mm,
        middle=1em,
        before skip=1.7em, after skip=1.2em,
        fontupper=\linespread{1.7}\selectfont
    }
}

\newtcolorbox{thmBox}[3][]{
  commonbox,
  colback=LogicBg,
  colbacktitle=LogicColor,
  before title={\refstepcounter{theorem}}, 
  title={#2 \thetheorem\ifstrempty{#3}{}{（#3）}},
  #1
}


\newcounter{theorem}[chapter]
\renewcommand{\thetheorem}{\thechapter.\arabic{theorem}}

\newenvironment{theorem}[1][]{\begin{thmBox}{定理}{#1}}{\end{thmBox}}
\newenvironment{lemma}[1][]{\begin{thmBox}{引理}{#1}}{\end{thmBox}}
\newenvironment{proposition}[1][]{\begin{thmBox}{命题}{#1}}{\end{thmBox}}
\newenvironment{corollary}[1][]{\begin{thmBox}{推论}{#1}}{\end{thmBox}}
\newenvironment{property}[1][]{\begin{thmBox}{性质}{#1}}{\end{thmBox}}
\newenvironment{induction}[1][]{\begin{thmBox}{归纳}{#1}}{\end{thmBox}}

\newtcolorbox{defBox}[3][]{
    commonbox, colback=DefBg, colbacktitle=DefColor,
    before title={\refstepcounter{theorem}},
  title={#2 \thetheorem\ifstrempty{#3}{}{（#3）}},
  #1
}

\newcounter{definition}[chapter]
\renewcommand{\thedefinition}{\thechapter.\arabic{definition}}

\newenvironment{definition}[1][]{\begin{defBox}{定义}{#1}}{\end{defBox}}
\newenvironment{axiom}[1][]{\begin{defBox}{公理}{#1}}{\end{defBox}}

\newtcolorbox{exBox}[3][]{
    commonbox, colback=ExBg, colbacktitle=ExColor,
    before title={\refstepcounter{theorem}},
    title={#2 \thetheorem\ifstrempty{#3}{}{（#3）}}, #1
}

\newtcolorbox{longexBox}[3][]{
    breakable, commonbox, colback=ExBg, colbacktitle=ExColor,
    before title={\refstepcounter{theorem}},
    title={#2 \thetheorem\ifstrempty{#3}{}{（#3）}}, #1
}

\newcounter{example}[chapter]
\renewcommand{\theexample}{\thechapter.\arabic{example}}

\newenvironment{example}[1][]{\begin{exBox}{例题}{#1}}{\end{exBox}}

\newenvironment{longexample}[1][]{\begin{longexBox}{例题}{#1}}{\end{longexBox}}


\newenvironment{exercise}[1][]{\begin{exBox}{练习}{#1}}{\end{exBox}}

\newtcolorbox{realexamBox}[2][]{
    commonbox, colback=ExBg, colbacktitle=ExColor,
    title={真题\ \ifstrempty{#2}{}{（#2）}}, #1
}

\newtcolorbox{longrealexamBox}[2][]{
    breakable, commonbox, colback=ExBg, colbacktitle=ExColor,
    title={真题\ \ifstrempty{#2}{}{（#2）}}, #1
}

\newenvironment{realexam}[1][]
{%\clearpage
\begin{realexamBox}{#1}}{\end{realexamBox}
%\clearpage
}

\newenvironment{longrealexam}[1][]
{%\clearpage
\begin{longrealexamBox}{#1}}{\end{longrealexamBox}
%\clearpage
}



\newenvironment{note}{
    \par\vspace{0.5em}\noindent
    \textbf{\color{NoteColor} 注：}
}{
    \par\vspace{0.5em}
}

\renewenvironment{proof}{
    \par\vspace{0.5em}\noindent
    \textbf{\color{LogicColor} 证明：}
}{
    \hfill \textcolor{LogicColor}{\rule{0.6em}{0.6em}} \par
}

\pagestyle{fancy}
\fancyhf{}
\fancyhead[L]{\small \kaishu \LectureNotesTitle}
\fancyhead[R]{\small \kaishu \AuthorInfo}
\fancyfoot[C]{\thepage}
\renewcommand{\headrulewidth}{0.5pt}
\renewcommand{\headrule}{\hbox to\headwidth{\color{LogicColor}\leaders\hrule height \headrulewidth\hfill}}

\fancypagestyle{plain}{
    \fancyhf{}
    \fancyhead[L]{\small \kaishu \LectureNotesTitle}
    \fancyhead[R]{\small \kaishu \AuthorInfo}
    \fancyfoot[C]{\thepage}
    \renewcommand{\headrulewidth}{0.5pt}
    \renewcommand{\headrule}{\hbox to\headwidth{\color{LogicColor}\leaders\hrule height \headrulewidth\hfill}}
}

\numberwithin{equation}{chapter}


% ======================================
%  ↑↑↑ 模板核心设置（可根据实际需求修改） ↑↑↑
% ======================================


\begin{document}

\frontmatter

% 显示标签页（可选，但建议添加，可以在代码最上方约第12行的地方开始修改设置）
\MakeTitlePage


% ====================
%  ↓↓↓ 前言（可选） ↓↓↓
% ====================

\chapter*{前言}
此笔记旨在复习考研数学一的同时记录下知识大纲，经典、易错的题目，常用的公式以及在学习过程中的总结，便于后续的再次复习。

为了节约时间，并没有必要将知识点的详细内容全盘描述，将知识框架搭建即可。

学习安排：先做基础题，标记遇到的生疏的知识点，然后使用讲义进行补充知识点并总结，然后做强化题目。

% 此模板系编者 Maki 为数学讲义设计之 \LaTeX\ 排版方案。经反复调试，现已基本定型，特以本文档集中展示其版式与功能，以便同好取用参考。

% 模板以 \texttt{ctexbook} 文档类为基础，采用 XeLaTeX 编译，正文使用 Noto Serif CJK SC 与 TeX Gyre Pagella 字体，章节标题以楷书排印。数学环境方面，定义、定理、例题、真题等均以 \texttt{tcolorbox} 实现彩色盒子排版，配色各异、层次分明；另设填空、选择等题型辅助命令，可直接用于试卷与讲义编排。讲义图片均以 TikZ 及 \texttt{pgfplots} 绘制，力求清晰可控。

% 本模板由 B 站账号「Maki 的完美算术教室」免费公开，允许商用；如若据此编写材料或二次开发，无需注明出处，可自由使用或二次开发。模板仍在持续迭代，若有建议或发现问题，欢迎来信告知（maki@maki-math.com）。

\vspace{3cm}

\hfill
\begin{minipage}{5cm}
    \centering
    sunlight \\
    二〇二六年五月\\
\end{minipage}

% ====================
%  ↑↑↑ 前言（可选） ↑↑↑
% ====================

% 目录（可选）
\tableofcontents

\mainmatter

% ===================================
%  ↓↓↓ 正文部分（开始你伟大的创作吧） ↓↓↓
% ===================================


% ===================================
\chapter{无穷级数}
% ===================================

本章分为：常数项级数、审敛方法、幂级数和收敛域、幂级数求和、函数展开为幂级数、傅里叶级数.

\section{常数项级数}
\begin{definition}
    [收敛]
    对于常数项级数 $\sum_{n=1}^{\infty} u_n$，如果部分和 $S_n$ 的极限$\lim_{n \to \infty} S_n$存在，则称该级数收敛；否则称该级数发散。
\end{definition}
\begin{induction}
    [收敛的性质]
    (1)收敛的线性组合收敛；
    
    (2)前有限项不影响敛散性；
    
    (3)收敛级数加括号后依旧收敛，但对于发散级数不成立($u_n = (-1)^n$)；
    
    (4)收敛的必要条件是：级数的项趋于零。
\end{induction}

\section{审敛方法}
\begin{induction}
    [正项级数的审敛方法]
    $\bigstar$在使用前必须要确保级数为正项级数。

    (1)充要条件：正项级数 $\sum_{n=1}^{\infty} u_n$ 收敛的充要条件是部分和 $S_n$ 有上界；
    
    (2)比较审敛法：大的收敛，小的也收敛；小的发散，大的也发散;
    
    极限形式：设 $\lim_{n \to \infty} \dfrac{u_n}{v_n} = l$，则：
    $0 < l < +\infty$ 时，$\sum u_n$ 与 $\sum v_n$ 同敛散；

    $l = 0$ 时，相当于$ u_n \le v_n$；
    $l = +\infty$ 时，相当于$ u_n \ge v_n$；

    (3)比值审敛法：设 $\lim_{n \to \infty} \dfrac{u_{n+1}}{u_n} = \rho $，则 $\rho < 1$ 收敛，$\rho > 1$ 发散，$\rho = 1$ 无结论；
    
    (4)根值审敛法：设 $\lim_{n \to \infty} \sqrt[n]{u_n} = \rho$，则 $\rho < 1$ 收敛，$\rho > 1$ 发散，$\rho = 1$ 无结论；
    
    (5)积分审敛法：
    
    设 $f(x)$ 在 $[1, +\infty)$ 上连续、正值、单调递减，则 $\sum_{n=1}^{\infty} f(n)$ 与 $\int_1^{\infty} f(x) dx$ 同敛散；
\end{induction}
\begin{note}
    比值审敛法的极限形式常用$0 < l < +\infty$进行同敛散的等价无穷小替换化简式子，例如：
    
    $e^{n^{-\alpha}\ln n}-1 ~\sim~ \ln (1 + n^{-\alpha}\ln n)~\sim~ \dfrac{\ln n}{n^{\alpha}}, n \to \infty$ ；
    
    根值审敛法的根号是$\sqrt[n]{u_n}$，而不是$\sqrt{u_n}。$
\end{note}

\begin{induction}
    [交错级数审敛法：莱布尼茨]
    交错级数$\sum(-1)^n u_n(u_n > 0)$ 收敛的充要条件是：

    (1) $ {u_n} $ 单调递减；
    
    (2) $\lim_{n \to \infty} u_n = 0$。
\end{induction}
\begin{note}
    判断$ {u_n} $递减的方法：
    
    (1)$f(n) = u_n$，则判断$f(x)$在$[A, +\infty)$上单调递减；

    (2)判断$u_{n+1} - u_n < 0$；
    
    (3)判断$\dfrac{u_{n+1}}{u_n} < 1$。
\end{note}

\begin{induction}
    [绝对收敛与条件收敛]
    (1) 如果级数 $\sum |u_n|$ 收敛，则称级数 $\sum u_n$ 绝对收敛；
    
    (2) 如果级数 $\sum u_n$ 收敛，但 $\sum |u_n|$ 发散，则称级数 $\sum u_n$ 条件收敛。
\end{induction}

\begin{induction}
    [泰勒展开在审敛中的应用]
    将级数逐项进行泰勒展开前几项，得到几个级数的和，根据这几个级数的敛散性来判断原级数的敛散性，当然如果出现两个发散级数就无法判定。
    
    例如：$\sum \ln(1+\frac{(-1)^n}{\sqrt{n}})$ = $\sum \left(\frac{(-1)^n}{\sqrt{n}} - \frac{1}{2n} + o(\frac{1}{n})\right)$，其中$\sum \frac{(-1)^n}{\sqrt{n}}$收敛，$\sum \frac{1}{2n}$发散，所以原级数发散。
\end{induction}

\section{幂级数和收敛半径}
\begin{definition}
    [幂级数]
    幂函数项级数简称幂级数，一般形式为$\sum_{n=0}^{\infty}a_n x^n$,或者$$\sum_{n=0}^{\infty}a_n (x-x_0)^n$$

    幂级数其实是一种级数形式的函数，其敛散性由于$x$的取值而不同，故幂级数的敛散性是一个函数问题.
    
    于是定义其收敛点为：使幂级数$\sum_{n=0}^{\infty}a_n (x-x_0)^n$收敛的$x$值；

    所有收敛点构成收敛区域；
\end{definition}

\begin{theorem}
    [阿贝尔定理]
    如果幂级数$\sum_{n=0}^{\infty}a_n x^n$在$x=x_1$处收敛，则在区间$|x|<|x-x_1|$内绝对收敛；
    
    如果在$x=x_2$处发散，则在区间$|x|>|x_2-x_0|$内发散。
\end{theorem}
\begin{note}
    阿贝尔定理的结论是：幂级数的收敛区间是一个以$x_0$为中心的区间，收敛区间的长度为$2R$，其中$R$为收敛半径。
\end{note}


\begin{induction}
    [一般级数$\sum u_n(x)$收敛区域的求法]
    (1)比值法：$\lim_{n \to \infty} \left|\dfrac{u_{n+1}(x)}{u_n(x)}\right| = \rho(x)$；
    
    (2)根值法：$\lim_{n \to \infty} \sqrt[n]{|u_n(x)|} = \rho(x)$。

    由$\rho(x)<1$初步得到收敛区域，而端点处(即为$\rho(x)=1$处)需要额外判定.
\end{induction}
\begin{note}
    对于幂级数$\sum a_n x^n$,$\quad \rho(x)$为$|x|$的一次式，于是特殊化为：
    
    (1)比值法：$\lim_{n \to \infty} \left|\dfrac{a_{n+1}}{a_n}\right| = \rho$；
    
    (2)根值法：$\lim_{n \to \infty} \sqrt[n]{|a_n|} = \rho$。

    收敛半径$R = \dfrac{1}{\rho}$,$\quad(-R, R)$的端点需要额外判定.
\end{note}

\begin{corollary}
    [幂级数变换与收敛的关系]
    已知幂级数$\sum a_n(x-x_1)^n$的收敛域,怎么得到$\sum b_n(x-x_2)^m$的收敛域呢？

    (1)平移不改变收敛半径，端点保持不变,例如：$$\sum a_n(x-x_1)^n \rightarrow \sum a_n((x-x_2))^n$$

    (2)称上或提出因子$(x-x_0)^k$不改变收敛半径，端点保持不变,例如：$$\sum a_n(x-x_0)^n \rightarrow \sum a_n(x-x_0)^{n+k} = (x-x_0)^k\sum a_n(x-x_0)^n$$

    (3)幂级数的求导与积分不改变收敛半径，端点需要额外判定,例如：$$\sum a_n(x-x_0)^n \rightarrow \sum na_n(x-x_0)^{n-1}, \quad \sum a_n(x-x_0)^n \rightarrow \sum \dfrac{a_n}{n+1}(x-x_0)^{n+1}$$
\end{corollary}


\section{幂级数求和}
\begin{definition}
    [级数的和函数]
    在收敛域上，记$S(x) = \sum_{n=1}^{\infty} u_n(x)，$称$S(x)$为级数的和函数;
\end{definition}

\begin{induction}
    [幂级数的运算方法]
    主要是对求和符号$\sum_{n=0}^{\infty}$的理解，$\sum$是一种运算规则，只对求和变量$n$有关，而$x$,常数均可视为求和无关常数；

    (1)提取因子：$\sum ka_n x^{n+k} = kx^k\sum a_n x^{n};$

    (2)逐项合并：$\sum a_n x^n + \sum b_n x^n = \sum (a_n + b_n)x^n$;

    (3)逐项相乘：$\sum a_n x^n \cdot \sum b_n x^n = \sum c_n x^n$，其中$c_n = a_0b_n + a_1b_{n-1} + ... + a_nb_0$；

    (4)逐项求导：$\dfrac{d}{dx}\sum a_n x^n = \sum na_n x^{n-1}$；

    (5)逐项积分：$\int_{0}^{x} \sum a_n t^n dt = \sum \dfrac{a_n}{n+1}x^{n+1}$;

    (6)变换下标：$\sum_{n=k}^{\infty} a_n x^n = \sum_{n=k+l}^{\infty} a_{n-l} x^{n-l}$;
\end{induction}

\begin{induction}
    [和函数的性质]
    (1)和函数在收敛域内连续；

    (2)和函数在收敛域内可导，且可导的结果仍为幂级数，\textbf{收敛半径不变}:
    $$\int_{0}^{x}S(t)dt = \int_{0}^{x}\sum_{n=0}^{\infty} a_n t^n dt = \sum_{n=0}^{\infty} \dfrac{a_n}{n+1}x^{n+1}$$

    (3)和函数在收敛域内可积，且可积的结果仍为幂级数，\textbf{收敛半径不变}:
    $$S'(x) = \dfrac{d}{dx}\sum a_n x^n = \sum_{n=1}^{\infty} na_n x^{n-1}$$
\end{induction}

\begin{induction}
    [幂级数求和的常用方法]
    求和函数$S(x) = \sum_{n=0}^{\infty} a_n x^n$，常用的方法有：
    
    (1)利用已知的幂级数求和公式，进行代换、积分、求导等运算,可以：
    
    i.先导后积：$S(x) = \int_{0}^{x} \left( \sum_{n=0}^{\infty} a_n t^n \right)' dt;$

    ii.先积后导：$S(x) = \left( \int_{0}^{x} \sum_{n=0}^{\infty} a_n t^n dt \right)';$

    (2)对$S(x) = \sum_{n=0}^{\infty} a_n x^n$求导，将右边化为$S(x),S'(x)$的形式，得到一个微分方程，解出$S(x)$；

    (3)对$S(x) = \sum_{n=0}^{\infty} a_n x^n$进行变换下标，得到一个递推公式，解出$S(x)$。
\end{induction}
\begin{note}
    由于积分、求导虽然不改变积分半径，但会改变端点的敛散性，所以在求和函数时不能只将用到的积分域求并，而是:
    
    应该记下积分半径，算出和函数$S(x)$之后再判断端点处$\sum a_n x^n$的敛散性。
\end{note}



\section{函数展开为幂级数}
\begin{definition}
    [函数展开为幂级数]
    如果函数$f(x)$在$x_0$处有任意阶导数，则$f(x)$在$x_0$处展开为幂级数的形式为：
    $$f(x) = \sum_{n=0}^{\infty} \dfrac{f^{(n)}(x_0)}{n!}(x-x_0)^n$$
    被称为$f(x)$在$x_0$处的\textbf{泰勒级数}。
    
    若$x_0 = 0$，则称为$f(x)$在$x_0$处的\textbf{麦克劳林级数}。
\end{definition}
\begin{note}
    此处的泰勒（麦克劳林）展开级数与微分中的泰勒（麦克劳林）展开式保持一致，区别在于：
    
    微分中的泰勒（麦克劳林）展开式是一个有限项的多项式，最后剩一个余项（无穷小或者$\exists \delta$形式来表示），
    而此处的泰勒（麦克劳林）展开级数是一个无穷项的级数。
\end{note}

\begin{induction}
    [展开方法]
    (1)公式直接法:直接套用泰勒（麦克劳林）展开公式，计算$f^{(n)}(x_0)$，然后代入公式即可；
    
(2)利用已知的幂级数展开式：通过代数运算、积分、求导等方法得到已知的函数，然后进行展开；
\end{induction}
\begin{note}
    $f(x) = \ln (1-x+x^2) = \ln(1+x^3) - \ln(1+x)$直接展开；
    
    $f(x) = \arctan\dfrac{1+x}{1-x} $则可以通过先求导、再展开、最后积分还原.
\end{note}


\section{傅里叶级数}
\begin{definition}
    [周期为$2l$的傅里叶级数]
    如果函数$f(x)$可以以$2l$为周期且在区间$[-l, l]$上可积，则$f(x)$可以展开为傅里叶级数：
    $$f(x) \backsim  \frac{a_0}{2} + \sum_{n=1}^{\infty} (a_n \cos \frac{n\pi x}{l} + b_n \sin \frac{n\pi x}{l}) = S(x)$$
    其中系数$a_n$和$b_n$由下列公式给出：
    $$a_n = \frac{1}{l} \int_{-l}^{l} f(x) \cos \frac{n\pi x}{l} dx, \quad b_n = \frac{1}{l} \int_{-l}^{l} f(x) \sin \frac{n\pi x}{l} dx$$
\end{definition}
\begin{note}
    不难发现：右边傅里叶级数的周期由最大的周期($a_1\cos\frac{\pi x}{l}, \quad b_1\sin\frac{\pi x}{l}$)决定,为$2l$，与$f(x)$周期一致;

    如果周期为$2l = 2\pi$，则傅里叶级数可以写为：
    $$f(x) \backsim  \frac{a_0}{2} + \sum_{n=1}^{\infty} (a_n \cos nx + b_n \sin nx)$$
    $$a_n = \frac{1}{\pi} \int_{-\pi}^{\pi} f(x) \cos nx dx, \quad b_n = \frac{1}{\pi} \int_{-\pi}^{\pi} f(x) \sin nx dx$$
\end{note}

\begin{theorem}
    [狄利克雷收敛定理]
    $f(x)$是以$2l$为周期且在区间$[-l, l]$上可积的函数，则$f(x)$的傅里叶级数$S(x)$处处收敛于$f(x)$的\textbf{平均值}：
    $$S(x) = \frac{f(x^+) + f(x^-)}{2}$$
\end{theorem}

\begin{induction}
    [函数的正弦展开、余弦展开]
    
    若函数$f(x)$定义在$[0,l]$,根据傅里叶级数的展开要求，被展开函数必须在$[-l,l]$上可积，为了将$f(x)$展开为正弦级数(余弦级数)，需要先进行周期奇(偶)延拓，然后再进行展开并将变量限制在$[0,l]$上.

    (1)周期奇延拓：
    $F(x) = \begin{cases}
        f(x), & 0 < x \le l,\\
        0, & x = 0,\\
        -f(-x), & -l \le x < 0,
    \end{cases}$
    $$a_n = \frac{1}{l} \int_{-l}^{l} F(x) \cos \frac{n\pi x}{l} dx = 0, \quad b_n = \frac{1}{l} \int_{-l}^{l} F(x) \sin \frac{n\pi x}{l} dx = \frac{2}{l} \int_{0}^{l} f(x) \sin \frac{n\pi x}{l} dx,$$
    $$f(x) \backsim \sum_{n=1}^{\infty} b_n \sin \frac{n\pi x}{l}.$$
    (2)周期偶延拓：
    $F(x) = \begin{cases}
        f(x), & 0 \le x \le l,\\
        f(-x), & -l \le x < 0.
    \end{cases}$
    $$a_n = \frac{1}{l} \int_{-l}^{l} F(x) \cos \frac{n\pi x}{l} dx = \frac{2}{l} \int_{0}^{l} f(x) \cos \frac{n\pi x}{l} dx, \quad b_n = \frac{1}{l} \int_{-l}^{l} F(x) \sin \frac{n\pi x}{l} dx = 0,$$
    $$f(x) \backsim \frac{a_0}{2} + \sum_{n=1}^{\infty} a_n \cos \frac{n\pi x}{l}.$$

\end{induction}










































\end{document}